\section{Introduction}

A black-box, LLM-guided grammar-fuzzing pipeline can be evaluated along two
different generalization axes. One is swapping the target implementation
while keeping everything else fixed -- a prior study did this, reusing an
identical pipeline unmodified against a second, independently implemented C
JSON parser, and found the mechanics of the loop (harness contract, proxy
signal, sandbox) transferred with no redesign~\priorworkcite. The
other axis, which that same study named explicitly as unaddressed future
work, is swapping the \emph{oracle} itself: every component of that
pipeline's failure signal -- a sanitizer-instrumented build, a subprocess's
exit code, a fatal signal -- assumes a native process that can crash in the
memory-unsafe sense. A memory-safe language has no such crash to look for,
and ``failure'' has to be redefined around something else entirely:
uncaught exceptions, silently violated API contracts, or values that
decode incorrectly without any exception at all.

This paper takes up that second axis directly, against a target chosen for
practical relevance rather than novelty for its own sake: JSON
deserialization in Dart, the language Flutter mobile and web applications
are written in. A first attempt at porting the original pipeline's crash
oracle onto a single, well-audited Dart JSON decoder would likely reproduce
the same kind of result the original study already reported for its second
C target -- a negative result, explainable but not informative, since a
heavily-used standard-library decoder is exactly the kind of target
unlikely to crash under a bounded fuzzing budget. We reject that port and
design two oracles instead that do not depend on a crash occurring at all.
The first is \emph{differential divergence}: rather than one implementation
whose accept/reject decision is the only signal available, we run every
generated document through four independently-implemented, widely-used Dart
JSON deserialization paths simultaneously -- hand-written parsing, and the
\texttt{json\_serializable}, \texttt{freezed}, and \texttt{built\_value}
codegen frameworks -- and treat disagreement between them, not a crash, as
the signal. The second is a \emph{ground-truth numeric round-trip check}:
JSON permits arbitrary-precision integer literals syntactically, and we
check whether a decoded value preserves the exact input magnitude, which
needs no assumption about what ``correct'' looks like beyond the literal
text the generator wrote.

Retargeting the oracle this way changes what the refinement loop is actually
being asked to optimize, and we treat that as an empirical question rather
than assuming the original result transfers. Concretely, this paper asks
three research questions. \textbf{RQ1}: reusing the original acceptance-rate
proxy and refinement loop essentially unmodified, does LLM-guided refinement
still raise a Dart JSON generator's rate of producing syntactically valid
documents, the way it did for the original C target? \textbf{RQ2}: across
the four Dart deserialization paths, does malformed or boundary-case input
surface behavioral divergence, and does LLM-guided refinement -- with its
prompt corrected to make divergence the explicit objective, not just an
incidental metric shown alongside others -- find more of it than a static,
schema-aware generator designed by hand after studying a small number of
cases? \textbf{RQ3}: does numeric decoding across these four paths preserve
exact integer values at magnitudes JSON's grammar permits but native
integer types do not, and how far does this generalize once measured at
scale rather than on one hand-picked probe?

We answer RQ1 affirmatively and RQ2 negatively, and the contrast between the
two is the paper's central finding: the same refinement mechanism, same
proposer model, same iteration budget, transfers cleanly to one objective
and not to the other. RQ3's answer is a concrete, quantified, and -- to our
knowledge -- previously undocumented defect pattern across two of the four
most common Dart serialization frameworks in real use. Concretely, this
paper contributes:

\begin{itemize}
  \item Two coverage-free, crash-free oracles -- cross-implementation
    differential divergence and ground-truth numeric round-trip checking --
    designed specifically for a memory-safe target where the original
    pipeline's sanitizer-based crash oracle does not apply, and a two-stage
    harness contract (a shared JSON-syntax gate, then four independent
    schema-level attempts) that makes the divergence oracle computable
    automatically from four otherwise-ordinary Dart programs.
  \item A budget-scoped, seeded, $n{=}5$-per-arm empirical comparison
    (RQ1) showing the original pipeline's refinement mechanism reproduces
    its acceptance-rate effect almost exactly on a new runtime and
    language -- 50.0\% to a 92.0\% mean, complete rank separation, exact
    two-sided $p \approx 0.0079$ -- a direct generalization result along the
    axis the original study left open.
  \item A second $n{=}5$-per-arm comparison (RQ2) showing the same
    mechanism, applied to a differential divergence objective instead,
    \emph{underperforms} a static, schema-aware generator by an order of
    magnitude (2.9\% mean vs.\ 22.2\% mean, complete rank separation in the
    opposite direction, exact two-sided $p \approx 0.0079$).
  \item A follow-up ablation that tests, rather than assumes, why: lifting
    the sandbox's restriction against importing a JSON library eliminates
    the LLM's generation-quality overhead almost entirely (every run reaches
    at or near 100\% syntactically valid output) yet leaves the divergence
    rate statistically unchanged (2.2\% mean, $p \approx 0.69$ against the
    constrained sandbox) -- establishing decisively that RQ2's shortfall is
    a targeting problem, not a generation-quality one.
  \item A second follow-up testing the natural next hypothesis -- that
    richer feedback (perturbation \emph{categories}, not just a score)
    closes the targeting gap -- and finding the opposite: a measurably
    \emph{worse} divergence rate (0.95\% mean, $p \approx 0.032$ against the
    score-only prompt), a genuine negative data point against the
    assumption that more informative feedback is strictly better.
  \item A methodological note, confirmed rather than assumed: an initial,
    unmodified refinement prompt found \emph{zero} divergence across three
    full seeded runs despite driving up unrelated rejection counts; a
    principled prompt revision that isolates divergence as an explicit
    score, validated on a held-out seed before any further spend, recovered
    a nonzero and consistent effect. Reporting the failed first design
    alongside the fix is itself part of this contribution.
  \item Two concrete, practitioner-actionable defects (RQ3 and part of
    RQ2), found by a static generator with no LLM involved and confirmed at
    scale by re-analyzing already-collected campaign data rather than
    running a redundant pipeline: \texttt{built\_value} silently defaults a
    missing list-typed required field instead of rejecting the input, and
    \texttt{json\_serializable}/\texttt{freezed} silently saturate an
    out-of-range integer field to \texttt{int64::MIN}/\texttt{MAX} with no
    exception, at a measured 5.7\% mismatch rate across 43{,}336 sampled
    checks spanning every nesting depth the schema's recursion produced.
  \item A released pipeline, harness, and full campaign artifact set,
    reusing the original project's sandboxing, seeding, and
    reproducibility-manifest machinery entirely unmodified where it is not
    C- or format-specific, and documenting precisely where it needed to
    change and where it did not.
\end{itemize}

The rest of the paper is organized as follows.
Section~\ref{sec:related-work} situates this work relative to coverage-guided,
grammar-based, differential, and LLM-assisted testing.
Section~\ref{sec:methodology} describes the oracle design, the four-path
harness, the generators, and what is reused unmodified from the original
pipeline. Section~\ref{sec:evaluation} reports all three research questions.
Section~\ref{sec:discussion} discusses the RQ1/RQ2 contrast, threats to
validity, and future work. Section~\ref{sec:conclusion} concludes.
